\chapimg{fig/ch24.jpg}
\chapterlead{\subhd{After this chapter you can} step the deterministic plant of layer P1 through a day with its equations P.1 to P.28, and rebuild the measurements of the virtual scanner of layer P2 (registration, voxels, volume, leaf area).}
\chapter{Phenotyping layers: the deterministic plant and the volumetric scan}\label{ch:pheno}
\markboth{Phenotyping}{Phenotyping}

Layer P1 grows a plant for $N=140$ days on a soil of eight layers of 15 cm, with no random numbers. Its equations are numbered P.1 to P.28 in the order the model uses them each day. Where the manual gives the form but not a species constant, the constant below is an \emph{assumed} typical value, and it is labelled.
\section{Development, light and biomass (P.1 to P.5)}

\exbox{Thermal time (P.1)}{Maize ($T_{\rm base}=8$, $T_{\rm opt}=30$ \textdegree C) grows under the (assumed) daily mean temperature $T(d)=18+9\sin(2\pi(d-60)/200)$. Find the degree days $d_g=\max(0,\min(T,T_{\rm opt}+6)-T_{\rm base})$ on days 0, 30, 70, 100 and the accumulated GDD at day 70 and day 139; if flowering needs $\mathrm{GDD}_{\rm flowering}=800$, on which day is it reached?}\sol  Day 0: $T=18+9\sin(-0.6\pi)=9.44$, $d_g=1.44$. Day 30: $T=10.72$, $d_g=2.72$. Day 70: $T=20.78$, $d_g=12.78$. Day 100: $T=26.56$, $d_g=18.56$. Summing: $\mathrm{GDD}(70)=346$ and $\mathrm{GDD}(139)=1536$ degree days. Flowering (800) is reached on day 99; the stage before flowering is $\mathrm{GDD}/800$, which is 0.43 on day 70.


\exbox{Reference evapotranspiration of the day (P.2)}{The reference evapotranspiration of the day (FAO-56, $G=0$) is $ET_0=\dfrac{0.408\Delta R_n+\gamma\frac{900}{T+273}u_2(e_s-e_a)}{\Delta+\gamma(1+0.34u_2)}$. For $T=28$ \textdegree C, $R_n=18$ MJ m$^{-2}$ d$^{-1}$, $u_2=2$ m s$^{-1}$, $\mathrm{RH}=45$ per cent and $\gamma=0.0674$ kPa \textdegree C$^{-1}$, compute it.}\sol  $e_s=0.6108\exp(17.27\times28/265.3)=3.780$ kPa; $e_a=0.45\,e_s=1.701$ kPa; $\Delta=4098\,e_s/(T+237.3)^2=0.2201$ kPa \textdegree C$^{-1}$. Radiation term: $0.408\times0.2201\times18=1.616$. Aerodynamic term: $0.0674\times900/301\times2\times(2.079)=0.838$. Denominator: $0.2201+0.0674(1+0.68)=0.3333$. So $ET_0=7.36$ mm d$^{-1}$ (about 66 per cent from radiation).


\exbox{Growth from light and from water (P.3 to P.5)}{Take $k=0.55$, RUE $=3.0$ g MJ$^{-1}$ PAR (assumed), $f_T=0.90$, $f_N=0.95$, $f_{\rm sal}=1$, PAR $=10$ MJ m$^{-2}$. Compute the light interception $f_i=1-e^{-k\,\mathrm{LAI}}$ and the radiation-limited growth $dW_{\rm rad}=\mathrm{RUE}f_Tf_Nf_{\rm sal}\mathrm{PAR}\,f_i$ for LAI $=0.5,\dots,5$. Then the water-limited growth $dW_{\rm wat}=(k_{te}/\mathrm{VPD})T_{\rm act}$ for $k_{te}=8$ Pa, VPD $=1.8$ kPa, $T_{\rm act}=3.5$ mm.}\sol  For LAI $=3$: $f_i=1-e^{-1.65}=0.808$, so $dW_{\rm rad}=3.0\times0.9\times0.95\times10\times0.808=20.7$ g m$^{-2}$ d$^{-1}$. Table~\ref{tab:rad} gives the others. For the water limit, $k_{te}/\mathrm{VPD}=8\ \mathrm{Pa}/1800\ \mathrm{Pa}=0.00444$ kg dry matter per kg water, and 3.5 mm of transpiration is 3.5 kg m$^{-2}$, so $dW_{\rm wat}=0.00444\times3.5=0.01556$ kg m$^{-2}$ $=15.6$ g m$^{-2}$. Growth is the smaller of the two, $dW=\min(dW_{\rm rad},dW_{\rm wat})$: for LAI $\le2$ light is limiting ($\le17.1$ g m$^{-2}$ d$^{-1}$ at LAI 2), for LAI $\ge3$ the 15.6 g of water limits it.\begin{table}[H]\centering\small\caption{Light interception and radiation-limited growth.}\label{tab:rad}
\begin{tabular}{lrr}\toprule
\hd{LAI} & \hd{$f_i$} & \hd{$dW_{\rm rad}$ (g m$^{-2}$ d$^{-1}$)}\\\midrule
0.5 & 0.240 & 6.17\\
1 & 0.423 & 10.85\\
2 & 0.667 & 17.11\\
3 & 0.808 & 20.72\\
4 & 0.889 & 22.81\\
5 & 0.936 & 24.01\\\bottomrule\end{tabular}\end{table}


\section{Soil water, leaf water and stomata (P.6 to P.10)}
\exbox{Leaf water potential and relative water content (P.8)}{$\psi=\pi_0/\mathrm{RWC}+P(\mathrm{RWC})$ with turgor $P(\mathrm{RWC})=-\pi_0(\mathrm{RWC}-\mathrm{RWC}_{\rm tlp})/(1-\mathrm{RWC}_{\rm tlp})$ above the turgor-loss point and $0$ below it. With $\pi_0=-1.2$ MPa and $\mathrm{RWC}_{\rm tlp}=0.80$ (assumed), tabulate $\psi$ against RWC and solve for the RWC at $\psi_{\rm leaf}=-1.5$ MPa. Also find $\psi_{\rm leaf}=\psi_{\rm root}-rT_{\rm act}/\mathrm{LAI}$ for $\psi_{\rm root}=-0.4$ MPa, $r=0.25$, $T_{\rm act}=4$ mm, LAI $=2$.}\sol  At RWC $=0.9$: $\pi=-1.2/0.9=-1.333$; $P=1.2(0.9-0.8)/0.2=0.6$; $\psi=-0.733$ MPa. Table~\ref{tab:rwc} gives the others. At $\psi=-1.5$ the bisection returns \textbf{RWC} $=0.800$ (exactly the turgor-loss point: there $P=0$ and $\psi=\pi_0/\mathrm{RWC}=-1.2/0.8=-1.5$ MPa). Flow through the plant: $\psi_{\rm leaf}=-0.4-0.25\times4/2=-0.9$ MPa, which corresponds to RWC $=0.878$.\begin{table}[H]\centering\small\caption{Pressure--volume curve of the leaf.}\label{tab:rwc}
\begin{tabular}{lrrr}\toprule
\hd{RWC} & \hd{$\pi_0/\mathrm{RWC}$} & \hd{$P$} & \hd{$\psi$ (MPa)}\\\midrule
1.00 & -1.200 & 1.200 & 0.000\\
0.95 & -1.263 & 0.900 & -0.363\\
0.90 & -1.333 & 0.600 & -0.733\\
0.85 & -1.412 & 0.300 & -1.112\\
0.80 & -1.500 & 0.000 & -1.500\\
0.75 & -1.600 & 0.000 & -1.600\\
0.70 & -1.714 & 0.000 & -1.714\\\bottomrule\end{tabular}\end{table}


\exbox{Stomata (P.9)}{The stomatal opening is $g_s=1/[1+e^{-b(\psi_{\rm leaf}-\psi_{50})}]$. With steepness $b=3$ MPa$^{-1}$ (assumed), compute $g_s$ for $\psi_{\rm leaf}=-0.5,\dots,-2.5$ and three genotypes: $\psi_{50}=-1.5$ (standard), $-1.0$ (a water saver that closes early) and $-2.0$ (a spender).}\sol  For the standard genotype at $-1.0$ MPa: $g_s=1/(1+e^{-3(-1.0+1.5)})=1/(1+e^{-1.5})=0.818$. At $\psi_{50}$ the opening is exactly 0.5. Table~\ref{tab:gs} lists all; at $-1.5$ MPa the water saver has $g_s=1/(1+e^{1.5})=0.182$ against $0.5$ for the standard genotype, so its transpiration demand is cut to $36$ per cent of the standard one at that potential, which saves water but also growth.\begin{table}[H]\centering\small\caption{Stomatal opening of three genotypes.}\label{tab:gs}
\begin{tabular}{lrrr}\toprule
\hd{$\psi_{\rm leaf}$} & \hd{$\psi_{50}=-1.5$} & \hd{$\psi_{50}=-1.0$} & \hd{$\psi_{50}=-2.0$}\\\midrule
-0.5 & 0.953 & 0.818 & 0.989\\
-1.0 & 0.818 & 0.500 & 0.953\\
-1.5 & 0.500 & 0.182 & 0.818\\
-2.0 & 0.182 & 0.047 & 0.500\\
-2.5 & 0.047 & 0.011 & 0.182\\\bottomrule\end{tabular}\end{table}


\exbox{Expansion growth (P.10)}{Expansion is $\frac1V\frac{dV}{dt}=\varphi\max(0,P-Y)$. With $\varphi=0.4$ d$^{-1}$ MPa$^{-1}$, yield threshold $Y=0.3$ MPa (assumed) and the turgor of the pressure--volume curve above (), find the relative growth rate and the leaf-area doubling time.}\sol  RWC $0.95$: $P=1.2(0.15)/0.2=0.9$, rate $=0.4(0.9-0.3)=0.24$ d$^{-1}$, doubling time $\ln2/0.24=2.89$ d. RWC $0.90$: $0.4(0.6-0.3)=0.12$ d$^{-1}$, $5.78$ d. RWC $0.85$: $P=0.3=Y$, rate $0$: growth stops, although the leaf still holds $85$ per cent of its water and the stomata are partly open. Because turgor falls before photosynthesis does, area and height slow first.


\exbox{Canopy temperature and CWSI (P.11)}{The canopy--air temperature difference is $T_c-T_a=\dfrac{r_aR_n}{\rho c_p}\dfrac{\gamma^*}{\Delta+\gamma^*}-\dfrac{\mathrm{VPD}}{\Delta+\gamma^*}$ with $\gamma^*=\gamma(1+r_c/r_a)$, and CWSI $=(dT-dT_{\rm low})/(dT_{\rm up}-dT_{\rm low})$. Take $T_a=30$ \textdegree C ($\Delta=0.2434$ kPa K$^{-1}$), $R_n=400$ W m$^{-2}$, $r_a=40$ s m$^{-1}$, $\rho c_p=1200$ J m$^{-3}$ K$^{-1}$, $\mathrm{VPD}=2.0$ kPa, $\gamma=0.0674$. Find $T_c-T_a$ for canopy resistances $r_c=0,50,100,200,500$ s m$^{-1}$ and the CWSI.}\sol  Limits: a transpiring canopy ($r_c=0$, $\gamma^*=\gamma$): $dT_{\rm low}=-3.54$ K; a non-transpiring one ($r_c\to\infty$, $\gamma^*\to\infty$): $dT_{\rm up}=r_aR_n/(\rho c_p)=13.33$ K. For $r_c=100$: $\gamma^*=0.0674\times3.5=0.2359$, so $dT=13.33\times0.2359/(0.2434+0.2359)-2.0/(0.2434+0.2359)=2.39$ K. Table~\ref{tab:cwsi} lists the others. The canopy temperature depression is $T_a-T_c=-dT$, positive for a canopy cooler than the air.\begin{table}[H]\centering\small\caption{Canopy temperature difference and stress index.}\label{tab:cwsi}
\begin{tabular}{lrr}\toprule
\hd{$r_c$ (s m$^{-1}$)} & \hd{$T_c-T_a$ (K)} & \hd{CWSI}\\\midrule
0 & -3.54 & 0.000\\
50 & 0.06 & 0.213\\
100 & 2.39 & 0.352\\
200 & 5.24 & 0.520\\
500 & 8.79 & 0.731\\\bottomrule\end{tabular}\end{table}


\section{Carbon, roots and soil strength (P.12 to P.14)}
\exbox{Carbon to the roots (P.12)}{The share of new carbon that goes to the roots is $f_r=\mathrm{clip}\bigl[(f_{r0}+(f_{rA}-f_{r0})\,\mathrm{stage})\,c_{\rm root}(1+0.9w_s),\,0.03,\,0.7\bigr]$ with $w_s=1-T_{\rm act}/T_p$. Wheat: $f_{r0}=0.42$, $f_{rA}=0.07$. Take $c_{\rm root}=1$ and find $f_r$ at stage $0,0.5,1$ for $T_{\rm act}/T_p=1,0.75,0.5$.}\sol  Stage $0.5$ and $T_{\rm act}/T_p=0.75$: base $0.42+(0.07-0.42)(0.5)=0.245$; stress $w_s=0.25$, factor $1+0.225=1.225$; so $f_r=0.245\times1.225=0.300$. The results are in Table~\ref{tab:fr}. Water shortage raises the root share (up to 45 per cent more at $T_{\rm act}/T_p=0.5$) and the stage lowers it: functional equilibrium.\begin{table}[H]\centering\small\caption{Carbon share to roots.}\label{tab:fr}
\begin{tabular}{lrrr}\toprule
\hd{Stage} & \hd{$T_{\rm act}/T_p$} & \hd{$w_s$} & \hd{$f_r$}\\\midrule
0 & 1.00 & 0.00 & 0.420\\
0 & 0.75 & 0.25 & 0.515\\
0 & 0.50 & 0.50 & 0.609\\
0.5 & 1.00 & 0.00 & 0.245\\
0.5 & 0.75 & 0.25 & 0.300\\
0.5 & 0.50 & 0.50 & 0.355\\
1.0 & 1.00 & 0.00 & 0.070\\
1.0 & 0.75 & 0.25 & 0.086\\
1.0 & 0.50 & 0.50 & 0.102\\\bottomrule\end{tabular}\end{table}


\exbox{Soil strength and root elongation (P.14)}{Soil strength is $Q=0.5\,e^{4.5(\rho_b-1.3)}[1+(1-S)]$ (MPa) and the relative elongation is $e/e_0=1/[1+(Q/2.5)^3]$. Compute both for bulk densities $1.2$ to $1.7$ g cm$^{-3}$ at relative water content $S=0.9$ and $0.6$, and find the bulk density that halves the elongation at $S=0.6$.}\sol  For $\rho_b=1.5$, $S=0.6$: $Q=0.5e^{0.9}(1.4)=0.5(2.4596)(1.4)=1.722$ MPa; $e/e_0=1/(1+(0.6887)^3)=1/(1.3267)=0.754$. Table~\ref{tab:Q} lists the rest. Halving requires $Q=2.5$: $0.5e^{4.5(\rho_b-1.3)}(1.4)=2.5$, $e^{4.5(\rho_b-1.3)}=3.571$, $\rho_b-1.3=\ln3.571/4.5=0.2831$, so $\rho_b=1.583$ g cm$^{-3}$. Drying the soil to $S=0.6$ from $0.9$ raises $Q$ by 27 per cent (factor $1.4/1.1$), as strength rises when the soil dries.\begin{table}[H]\centering\small\caption{Soil strength and root elongation.}\label{tab:Q}
\begin{tabular}{lrrr}\toprule
\hd{$\rho_b$} & \hd{$S$} & \hd{$Q$ (MPa)} & \hd{$e/e_0$}\\\midrule
1.2 & 0.9 & 0.351 & 0.997\\
1.2 & 0.6 & 0.446 & 0.994\\
1.3 & 0.9 & 0.550 & 0.989\\
1.3 & 0.6 & 0.700 & 0.979\\
1.4 & 0.9 & 0.863 & 0.961\\
1.4 & 0.6 & 1.098 & 0.922\\
1.5 & 0.9 & 1.353 & 0.863\\
1.5 & 0.6 & 1.722 & 0.754\\
1.6 & 0.9 & 2.122 & 0.621\\
1.6 & 0.6 & 2.700 & 0.442\\
1.7 & 0.9 & 3.327 & 0.298\\
1.7 & 0.6 & 4.235 & 0.171\\\bottomrule\end{tabular}\end{table}


\exbox{Root elongation (P.13)}{Root tip elongation is $\dfrac{dl}{dt}=e_0\,2^{(T_s-20)/10}\Bigl(\dfrac e{e_0}\Bigr)f_wf_{\rm age}$ with $f_w=\mathrm{clip}\bigl((\theta_{\rm rel}-0.06)/0.5;0.05,1\bigr)$. For an axis with $e_0=2.5$ cm d$^{-1}$ at $T_s=16$ \textdegree C, $\theta_{\rm rel}=0.5$, $e/e_0=0.754$ and $f_{\rm age}=0.8$, find the elongation and the time to reach 60 cm.}\sol  Temperature factor $2^{-0.4}=0.758$; $f_w=(0.5-0.06)/0.5=0.88$. Elongation: $2.5\times0.758\times0.754\times0.88\times0.8=1.006$ cm d$^{-1}$, so a tip needs about 60 days to reach 60 cm at these conditions; in the wetter, looser soil ($e/e_0=1$, $f_w=1$, $T_s=20$) the same axis would need $60/(2.5\times0.8)=30$ days.


\section{Irrigation, mulch, nitrogen and salt in the plant (P.15 to P.19)}
\exbox{Project irrigation event (P.16)}{Irrigation (P.16): if the sensor suction $h_s\ge h_{\rm trig}$ and the minimum interval has passed, the root zone is refilled to field capacity, $d_{\rm net}=\min(\sum(\theta_{fc}-\theta_z)\Delta z,D_{\max})$, and the gross depth is $I=d_{\rm net}/E_a$. The sprinkler has minimum interval 2 d, $D_{\max}=35$ mm, $E_a=0.75$; surface 3 d, $60$ mm, $0.60$; drip 1 d, 10 mm, $0.90$. Eight layers of $\Delta z=15$ cm have deficits $\theta_{fc}-\theta_z=0.02,0.03,0.04,0.05,0.05,0.04,0.03,0.02$. Find the net and gross depth for each method.}\sol  The total deficit is $\sum(\theta_{fc}-\theta_z)\Delta z=15(0.02+0.03+0.04+0.05+0.05+0.04+0.03+0.02)=15(0.28)=4.2$ cm $=42$ mm. Sprinkler: $d_{\rm net}=\min(42,35)=35$ mm, $I=35/0.75=46.7$ mm. Surface: $\min(42,60)=42$ mm, $I=42/0.60=70$ mm. Drip: $\min(42,10)=10$ mm, $I=10/0.90=11.1$ mm per event (and the drip line waters only when the layer at the emitter depth is at least 12 per cent depleted). The surface method needs 70 mm per event against 46.7 for the sprinkler for the same deficit; a drip system needs about four small events to replace the same deficit.


\exbox{Mulch (P.17)}{Mulch (P.17) gives $E_p'=f_mE_p$ and $P'=\max(0,P-C_m)$. For $f_m=0.4$, $C_m=2$ mm and a day with $E_p=4.5$ mm and a rain event of $P=1.5$, $4$ and $12$ mm, find the soil evaporation and the effective rain.}\sol  $E_p'=0.4\times4.5=1.8$ mm. $P'=\max(0,1.5-2)=0$ for the 1.5 mm event (the mulch keeps all of it), $2.0$ mm for the 4 mm event, $10.0$ mm for the 12 mm. Over a season with many small events the interception can be substantial: if the rain of 90 events of mean depth 1.5 mm (135 mm) never exceeds $C_m$, none of it reaches the soil.


\exbox{Salt in the root zone (P.19)}{The root-zone solution of volume $V$ (twice the plant-available water of the column) follows $V\,dC/dt=I\,EC_w+0.05P-DC$, $EC_e=C/2$, $f_{\rm sal}=1-\frac s{100}\max(0,EC_e-EC_{\rm th})$, $\psi_{\rm soil}\leftarrow\psi_{\rm soil}-0.036EC_e$. With $V=200$ mm (assumed), a daily irrigation of $I=4$ mm of $EC_w=2$ dS m$^{-1}$, no rain, a drainage of $D=1.5$ mm d$^{-1}$ and an initial $C=1$, integrate by daily steps and find the steady state; for wheat ($EC_{\rm th}=6$, $s=7.1$) compute $f_{\rm sal}$.}\sol  One step: $C_1=1+(4\times2-1.5\times1)/200=1.0325$. The steady state sets the right-hand side to zero: $C^*=I\,EC_w/D=8/1.5=5.33$ dS m$^{-1}$, so $EC_e=2.67$. The relaxation time is $V/D=200/1.5=133$ days, so after 60 days $C=2.57$ (the analytic $C(t)=C^*-(C^*-C_0)e^{-Dt/V}=1.31$ agrees). At $C^*$, $EC_e=2.67$ is below the threshold 6, so $f_{\rm sal}=1$ and the osmotic term is $-0.036\times2.67=-0.096$ MPa. If the drainage is cut to $0.3$ mm d$^{-1}$: $C^*=26.7$, $EC_e=13.3$, $f_{\rm sal}=1-0.071(13.3-6)=0.48$, a loss of 52 per cent.

